\subsection{Các Stakeholder và User story}
\subsubsection{Stakeholder}
Hệ thống bao gồm stakeholder là người chơi, chính là những người tham gia trò chơi, có thể là một hoặc nhiều người tùy theo số lượng người vào cùng 1 phòng trong trò chơi.
\subsubsection{User story của người chơi}
Ta có thể có được các user story đối với các chức năng, trường hợp trong trò chơi. Cụ thể về các user story:

\begin{itemize}
    \item  Đối với user story về các luồng thực thi ở trang chủ của trò chơi, là người chơi, tôi có thể:
        \begin{itemize}
            \item Chọn tạo phòng hoặc tham gia phòng của người chơi khác.
            \item Điều chỉnh các thông số cấu hình của trò chơi khi truy cập vào cài đặt.
            \item Xem thông tin cá nhân của người chơi.
            \item Xem thông tin về nhà phát triển của trò chơi.
        \end{itemize}
    \item Đối với user story về các luồng thực thi khi đã vào phòng của trò chơi, là người chơi, tôi có thể:
        \begin{itemize}
            \item Xem thông tin phòng chơi.
            \item Sử dụng các nhạc cụ có trong phòng.
            \item Thay đổi trang phục.
            \item Ngồi vào bàn và bắt đầu chơi (nếu đã đủ số lượng người chơi tối thiểu).
            \item Nếu là chủ phòng thì có thể kick người chơi khác khỏi phòng của mình.
        \end{itemize}
    \item Đối với user story về gameplay của trò chơi, là người chơi, tôi có thể:
        \begin{itemize}
            \item Xem hướng dẫn về trò chơi như luật, cách chơi về board game.
            \item Kết thúc màn chơi hiện tại. Sau khi kết thúc, tôi có thể chọn một trong các lựa chọn: bắt đầu lại hoặc rời khỏi ghế và làm những thứ khác.
        \end{itemize}
    \item Đối với user story về các tùy chỉnh cấu hình khi truy cập vào cài đặt trong trò chơi, là người chơi, tôi có thể:
        \begin{itemize}
            \item Điều chỉnh các thông số chung của trò chơi như âm nhạc, âm thanh.
        \end{itemize}
\end{itemize}